Задача 1. Из утверждений <<число а делится на 2>>, «число а делится на 4»,<<число а делится на 12>> и <<число а делится на 24>> три верных, а одно неверное. Какое?

\textbf{Решение:}
Пусть последнее утверждение верное. Но если $a$ делится на 24, то оно делится и на 12, и на 4 и на 2, что нам не подходит. Значит единственный возможный случай, когда это утверждение неверное. Такое может быть, например, если $a = 12.$

Задача 2. Как вы думаете, среди четырех последовательных натуральных чисел будет ли хотя бы одно делиться на 2? А на 3? А на 4? А на 5?

Задача 3. Известно, что р $>3$ и р - простое число, т.е. оно делится только на единицу и на себя само. Как вы думаете:
а) будут ли чётными числа ( $\mathrm{p}+1$ ) и ( $\mathrm{p}-1$ );
б) будет ли хотя бы одно из них делиться на 3;
в) будет ли хотя бы одно из них делиться на 4;
г) будет ли хотя бы одно из них делиться на 5?

Задача 4. Найдите все числа, при делении которых на 7 в частном получится то же число, что и в остатке.

Задача 5. Коля, Серёжа и Ваня регулярно ходили в кинотеатр. Коля бывал в нём каждый 3-й день, Серёжа - каждый 7-й, Ваня - каждый 5-й. Сегодня все ребята были в кино. Когда все трое встретятся в кинотеатре в следующий раз?

Задача 6. При делении некоторого числа m на 13 и 15 получили одинаковые частные, но первое деление было с остатком 8, а второе без остатка. Найдите число m.

Задача 7. В комнате стоят трехногие табуретки и четырехногие стулья. Когда на все эти сидячие места уселись люди, в комнате оказалось 39 ног. Сколько в комнате табуреток?

Задача 8. Ковбой Билл зашёл в бар и попросил у бармена бутылку виски за 3 доллара и шесть коробков непромокаемых спичек, цену которых он не знал. Бармен потребовал с него 11 долларов 80 центов (1 доллар= 100 центов), и в ответ на это Билл вытащил револьвер. Тогда бармен пересчитал стоимость покушки и исправил ошибку. Как Билл догадался, что бармен пытался его обсчитать?